% !TeX root = ../main.tex
\section{The app's pipeline end to end}
\label{sec:app}

This section puts the stages of Part~II together as they run in the app.
Figure~\ref{fig:app-pipeline} is the data flow; Table~\ref{tab:app-components}
lists each component, its file and what happens when an optional part is
missing.

\begin{figure}[h]
\centering
\begin{tikzpicture}[
  node distance=3.5mm,
  box/.style={draw, rounded corners=2pt, font=\small, align=center,
              minimum height=8mm, text width=40mm, inner sep=3pt, fill=white},
  opt/.style={box, dashed},
  arr/.style={-{Stealth[length=2mm]}, semithick}]
  \node[box] (in) {Audio file or YouTube\\(\code{yt-dlp})};
  \node[box, below=of in] (dec) {Decode + SHA-256};
  \node[opt, below=of dec] (sep) {Separation\\BS-RoFormer / Demucs};
  \node[opt, right=12mm of sep] (beat) {Beats, bars, meter\\Beat This! / DSP};
  \node[box, below=of sep] (pan) {Pan split\\(left / centre / right)};
  \node[box, below=of pan] (notes) {Guitar notes per part\\GAPS / v2 / Basic Pitch};
  \node[box, below=of notes] (tech) {Techniques\\(Basic Pitch contour)};
  \node[box, left=12mm of notes] (bass) {Bass notes\\Basic Pitch};
  \node[box, below=of tech] (arr) {Arrange: quantise,\\fingering, tuning, notation};
  \node[box, below=of arr] (out) {GP5 tab + \code{.gvtab.json}\\in the tab repository};
  \node[opt, right=12mm of out] (tone) {Tone estimate per guitar\\(in progress, optional)};
  \draw[arr] (in) -- (dec);
  \draw[arr] (dec) -- (sep);
  \draw[arr] (dec.east) -| (beat.north);
  \draw[arr] (sep) -- (pan);
  \draw[arr] (pan) -- (notes);
  \draw[arr] (notes) -- (tech);
  \draw[arr] (tech) -- (arr);
  \draw[arr] (sep.west) -| (bass.north);
  \draw[arr] (bass.south) |- (arr.west);
  \draw[arr] (beat.south) |- (arr.east);
  \draw[arr] (arr) -- (out);
  \draw[arr] (out) -- (tone);
\end{tikzpicture}
\caption{Data flow of \code{transcribe::run} (\file{transcribe/mod.rs}). Dashed
boxes are optional: the pipeline still runs without them, with the fallbacks of
Table~\ref{tab:app-components}.}
\label{fig:app-pipeline}
\end{figure}

\begin{table}[h]
\centering
\caption{Components of the transcription pipeline.}
\label{tab:app-components}
\small
\begin{tabularx}{\textwidth}{@{}>{\raggedright}p{2.2cm}>{\raggedright}p{3.6cm}>{\raggedright}p{3.6cm}X@{}}
\toprule
Stage & Default & File & Without the optional part \\
\midrule
Source & file (mp3, flac, ogg, wav, m4a) or YouTube & \file{transcribe/youtube.rs} & -- \\
Separation & BS-RoFormer SW (MSST, Python) & \file{transcribe/separation.rs} & Demucs \code{htdemucs_6s}; else the full mix, one guitar track \\
Guitar parts & pan split & \file{transcribe/stereo.rs} & one guitar track (mono input or no separation) \\
Beats and bars & Beat This! final0 (small0) & \file{transcribe/beat_net.rs} & DSP tracker, 4/4 only (\file{beats.rs}) \\
Guitar notes & GAPS CRNN; v2 if chosen & \file{transcribe/guitar_crnn.rs} & Basic Pitch \\
Bass notes & Basic Pitch & \file{transcribe/basic_pitch.rs} & -- \\
Techniques & contour rules & \file{transcribe/techniques.rs} & -- \\
Fingering & learned tagger (6 strings) & \file{transcribe/fingering_net.rs} & Viterbi (\file{arrange.rs}) \\
Notation & quantise, tuplets, ties, tuning & \file{transcribe/arrange.rs} & -- \\
Output & GP5 + sidecar & \file{shared/tab/export.rs}, \file{repository.rs} & -- \\
\bottomrule
\end{tabularx}
\end{table}

\subsection{Entry points}

\begin{itemize}
\item \textbf{Tabs page}, button ``Transcrever áudio'', and the Spotify page
  for a track without a tab (\file{frontend/ui/transcribe.rs}). The form lets
  the user pick a file or search YouTube, and choose the guitar model and the
  fingering engine.
\item \textbf{Batch CLI}, which skips audio already in the repository (also
  \code{--youtube}):
\begin{verbatim}
cargo run --release -p guitar_desktop --bin transcribe -- <paths>
\end{verbatim}
\item \textbf{YouTube} (\file{transcribe/youtube.rs}): wraps the external
  \code{yt-dlp}; searches without an API key, ranks results by title words,
  ``Artist -- Topic'' channels, penalties for live/cover/karaoke and closeness
  to the Spotify duration; downloads audio into a private cache, never into
  the repository. Video IDs are validated before reaching the command line.
\end{itemize}

\subsection{What a transcription produces}

The result goes to the local tab repository
(\file{<data dir>/guitar_virtual_amp/tab_repository/}):
\begin{itemize}
\item \file{Artist - Title.gp5}, a Guitar Pro 5 file written through the
  \code{guitarpro} crate (opens in Guitar Pro, TuxGuitar, MuseScore). Tempos are
  rounded per bar, as GP5 requires.
\item \file{Artist - Title.gvtab.json}, the \term{sidecar} (\code{TabMetadata}):
  the source audio's SHA-256 (to skip or replace repeats), the engines used
  (separator, note model, beat model, fingering), the YouTube ID when there is
  one, and the time in the audio of every bar (\code{bar_audio_secs}). The tab
  player uses the bar times (\code{TabMetadata::sync_map}) to play the original
  recording in exact sync, whatever the GP5 tempo rounding.
\end{itemize}

\subsection{Cost}

On the laptop (RTX 4070, 8\,GB): a 4.4-minute song takes about 70\,s and
2.4\,GB of RAM \measured; BS-RoFormer separation alone is 34\,s on a 261\,s
song. The CRNN windows are the memory peak (about 600\,MB each in tract, so at
most 4 at a time). The tract crates are built with \cmd{opt-level = 3} even in
debug builds, otherwise transcription takes minutes.

\begin{inapp}[: tone estimate, experimental \done]
An optional last stage: with
``Estimar o timbre (experimental)'' checked in the Transcribe form (setting
\code{AppSettings::transcription_tone}), each guitar part's audio and notes go
to \code{tone::estimate_from_transcription}, and the estimated chain is saved as
a user preset tagged with the song and the track (\code{tone::save_preset}), so
the Tabs page can suggest it. It needs separation. The method is the subject of
the companion document \file{docs/papers/tone-estimation/}.
\end{inapp}
